\documentclass[review]{elsarticle}
\usepackage{booktabs}
\usepackage{longtable}
\usepackage{graphicx}
\usepackage{hyperref}
\usepackage{array}
\usepackage{multirow}

\journal{Journal of Biomedical Informatics}

\begin{document}

\begin{frontmatter}

\title{Factuality Preservation in Patient-Facing Medical Text Simplification: A Systematic Methodological Review with a Low-Resource Bangla Evidence Map}

\author[inst1]{First Author}
\author[inst1]{Second Author}
\address[inst1]{Bangla Clinical NLP Initiative}

\begin{abstract}
\textbf{Objective:} To systematically review methods, datasets, and evaluation practices for factuality-preserving medical text simplification and assess transferability to Bangla healthcare NLP. 
\textbf{Methods:} We will follow PRISMA 2020 and PRISMA-S. Searches cover ACL Anthology, PubMed/PMC, Scopus, ScienceDirect, ACM Digital Library, and IEEE Xplore for 2020--2025 peer-reviewed papers. Preprints and short weak conference papers are excluded from the analytical corpus. 
\textbf{Results:} The planned analytical corpus contains 80 selected peer-reviewed studies grouped by clinical simplification, factuality evaluation, medical NLI, Bangla/Indic healthcare NLP, RAG/KG grounding, constrained decoding, datasets, and evaluation methods. 
\textbf{Conclusion:} The review will identify evidence gaps in relation-level factuality, Bangla clinical simplification resources, and human evaluation rigor.
\end{abstract}

\begin{keyword}
medical text simplification \sep factuality \sep clinical NLP \sep Bangla NLP \sep low-resource languages \sep systematic review
\end{keyword}

\end{frontmatter}

\section{Introduction}
Clinical documents are difficult for patients to understand, and factual errors in simplification can introduce clinical risk. This review focuses on factuality preservation across entities, relations, and attributes, with a low-resource Bangla evidence map.

\section{Methods}
\subsection{Protocol and reporting}
This review follows PRISMA 2020 and PRISMA-S. The search strategy and full screening log are provided in Supplementary Files S1--S4.

\subsection{Eligibility criteria}
The analytical corpus includes peer-reviewed journal articles and full-length conference papers from 2020--2025. Preprint-only papers, editorials, posters, non-archival demos, and weak short conference papers are excluded.

\subsection{Search strategy}
The search combines clinical/patient-facing text, simplification/explanation, factuality/faithfulness, and methods. Full Boolean strings are in Supplementary File S1.

\subsection{Screening and data extraction}
Two reviewers independently screen titles/abstracts and full texts. The extraction codebook includes citation, venue, task, architecture, dataset, language, factuality method, readability metric, human evaluation, risk of bias, and Bangla relevance.

\subsection{Risk of bias and certainty}
Benchmark/model papers are assessed using a custom NLP-Bias tool. GRADE certainty is rated for factuality, readability, patient comprehension, Bangla transferability, and reproducibility.

\section{Results}
\subsection{Search flow}
Insert PRISMA Figure 1 here after the locked database search.

\subsection{Characteristics of the included corpus}
Use Table 2 from \texttt{top80\_selected\_corpus.csv}.

\subsection{Factuality preservation taxonomy}
Use Figure 2 and the factuality matrix from the workbook.

\subsection{Dataset and resource landscape}
Use Table 3 from \texttt{dataset\_inventory.csv}.

\subsection{Model and method comparison}
Use Table 4 from \texttt{model\_comparison\_table.csv}.

\section{Discussion}
Discuss gaps G1--G7 and their relationship to future Bangla clinical simplification systems.

\section{Conclusion}
This review will provide a reproducible evidence map for factuality-preserving clinical text simplification and a concrete research agenda for low-resource Bangla healthcare NLP.

\bibliographystyle{elsarticle-num}
\bibliography{selected_top80_corpus}

\end{document}
